\chapter{Конформная мода и Majorana-rank: тест критерия Ф.3}

После отрицательного phase-to-mass gate остаётся проверить вторую часть
предложенной 4D-интерпретации: может ли удаление одной conformal/gauge mode
превратить полный Majorana-модуль размерности \(24\) в физический rank \(23\)?

\section{Diffeomorphism gauge complex}

В четырёх измерениях metric fluctuation имеет десять локальных компонент,
а линейное действие diffeomorphism задаётся
\begin{equation}
 \delta_\xi h_{\mu\nu}=\nabla_\mu\xi_\nu+\nabla_\nu\xi_\mu.
\end{equation}
На Fourier-level с momentum \(k\) его symbol равен
\begin{equation}
 G_k(\xi)_{\mu\nu}=k_\mu\xi_\nu+k_\nu\xi_\mu.
\end{equation}
Для \(k\neq0\) этот map имеет rank \(4\), но pure trace tensor
\(h_{\mu\nu}=c g_{\mu\nu}\) не принадлежит его image. Для constant mode
\(k=0\) symbol вообще равен нулю. Глобально тот же факт следует из
\begin{equation}
 \delta_\xi\Vol(K)
 =\int_K\nabla_\mu\xi^\mu\,d\Vol=0
\end{equation}
на компактном manifold без границы. Следовательно, constant volume mode не
является diffeomorphism gauge direction.

\section{Почему gauge volume не действует на \texorpdfstring{\(\mathbb R^{24}\)}{R24}}

Минимальное поле конфигураций имеет direct-sum структуру
\begin{equation}
 \operatorname{Sym}^2(T^*K)\oplus\mathcal M_{Maj},
 \qquad \dim_{\mathbb R}\mathcal M_{Maj}=24.
\end{equation}
В замороженной постановке BRST/diffeomorphism map действует в первом блоке:
\begin{equation}
 \xi\longmapsto(G\xi,0).
\end{equation}
Поэтому quotient по \(\operatorname{Diff}(K)\), ghost determinant и
gauge-volume factor не меняют rank второго summand:
\begin{equation}
 \dim\mathcal M_{Maj}^{physical}=24,
\end{equation}
если отсутствует отдельный constraint или projector внутри самого
Majorana-модуля.

\section{Weyl и wrong-sign не спасают подсчёт}

Pure trace становится gauge direction только при дополнительной локальной
Weyl symmetry
\begin{equation}
 \delta_\sigma g_{\mu\nu}=2\sigma g_{\mu\nu}.
\end{equation}
Такая symmetry отсутствует в текущем massive parent action: фиксированные
радиусы, Majorana/Yukawa scales и dimensionful defect couplings нарушают её.
Даже если формально добавить Weyl quotient, он rescale-ит поля по их Weyl
weights, но не выбирает одну вещественную polarization внутри
\(\mathbb R^{24}\). На zero-field background линейный gauge map в
Majorana-блок равен нулю.

Также «неправильный знак» Euclidean conformal factor означает выбор или
rotation integration contour. Это не gauge quotient и не уменьшение
конечномерного rank.

Наконец, коэффициент \(1/24\) в
\begin{equation}
 \widehat A=1-\frac{p_1}{24}+\cdots
\end{equation}
является нормировкой characteristic class, а не числом Majorana degrees of
freedom. Из него не следует операция \(24-1\).

\begin{theorem}[Conformal-rank no-go]
В 4D parent architecture без точной локальной Weyl gauge symmetry
обработка conformal factor и деление на gravitational gauge volume не
превращают \(\mathcal M_{Maj}\simeq\mathbb R^{24}\) в rank-\(23\) модуль.
Для такого сокращения необходим rank-one operator, действующий внутри
Majorana fiber.
\end{theorem}

\section{Сравнение с defect-route}

Предыдущий dimension gate уже доказал, что удаление generation singlet
вычитает полный внутренний spinor module, а не одну вещественную компоненту.
Но class-D defect на систолической геодезической даёт условный обход этого
no-go: transverse mod-two kernel, longitudinal periodic real mode,
generation singlet и self-dual cycle line имеют combined rank \(1\). Внутри
этой defect-модели дополнение в \(\mathbb R^{24}\) действительно имеет rank
\(23\).

Различие принципиально. Conformal factor находится в metric field space и
не действует на Majorana fiber. Defect zero mode является собственным
вектором Majorana/BdG operator и потому способен определить настоящий
rank-one projector. Его оставшийся gap --- вывод odd-winding mass order
parameter и restriction map из общего S2T action.

\begin{protocol}[Условие переоткрытия Ф.3]
Конформная ветвь может быть переоткрыта только если:
\begin{enumerate}
 \item полное massive действие обладает точной Weyl gauge symmetry;
 \item построен её BRST-complex вместе с diffeomorphism ghosts;
 \item mixed gauge map имеет rank-one image внутри \(\mathcal M_{Maj}\);
 \item quotient удаляет ровно одну реальную Majorana direction, а не общий
 scale или полный internal multiplet;
 \item тот же complex независимо фиксирует relative trace--Hodge weight.
\end{enumerate}
Альтернативный допустимый путь --- вывести уже построенный class-D defect из
parent action без нового выбранного order parameter.
\end{protocol}

Текущий вердикт Ф.3 отрицателен для conformal/gauge-volume объяснения.
Число \(23\) не выводится из \(\widehat A\)-знаменателя и не возникает при
удалении gravitational conformal mode. Оно сохраняет условный статус только
через явный rank-one Majorana defect либо иной operator в \(\mathbb R^{24}\).

Машинный протокол сохранён в
\texttt{s2t\_conformal\_majorana\_rank\_gate\_results.json}.